% 原书物理页=印刷页63--68；2.5 全节.
% 注：2.5.7 之命题的证明末尾溢出至原书第 69 页顶部（2.6 节标题之前），一并译出.
\section{半单 Lie 代数的结构}\label{sec:lie-2-5}

在这一部分中，$K$ 表示一个特征为零的代数闭域.

\begin{numpara}[Cartan 子代数]\label{np:lie-2-5-1}
设 $\mathfrak g$ 是一个 $K$-Lie 代数. 若 $\mathfrak h\subset\mathfrak g$ 是一个 Lie 子代数，则记
$N(\mathfrak h):=\{X\in\mathfrak g,\ \operatorname{ad}_{X}(\mathfrak h)\subset\mathfrak h\}$
为它的\newterm{lie-normalizer}{正规化子}{normalizer}，它也是 $\mathfrak g$ 的一个 Lie 子代数（请自行验证）.
\end{numpara}

\begin{definition*}
$\mathfrak g$ 的一个 \newterm{lie-cartan-subalgebra}{Cartan 子代数}{Cartan subalgebra} 是一个幂零 Lie 子代数 $\mathfrak h$，满足 $N(\mathfrak h)=\mathfrak h$.
\end{definition*}

我们将证明这样的子代数确实存在. 对 $X\in\mathfrak g$，记
$\mathfrak g_{X,0}:=\bigcap_{n\in\mathbb N}\operatorname{Ker}(\operatorname{ad}_{X})^{n}\subset\mathfrak g$
为 $\operatorname{ad}_{X}$ 属于特征值 $0$ 的广义特征空间，再令
\[
\operatorname{rk}(\mathfrak g):=\min_{X\in\mathfrak g}\bigl(\dim_{K}(\mathfrak g_{X,0})\bigr).
\]
注意 $\dim_{K}(\mathfrak g_{X,0})$ 是 $0$ 在 $\operatorname{ad}_{X}$ 的特征多项式中的重数. 整数 $\operatorname{rk}(\mathfrak g)$ 称为 $\mathfrak g$ 的\newterm{lie-rank}{秩}{rank}，而元素 $X$ 称为\newterm{lie-regular-element}{正则}{regular}的，若 $\operatorname{rk}(\mathfrak g)=\dim_{K}(\mathfrak g_{X,0})$.
\begin{exercise*}
证明 $\operatorname{rk}(\mathfrak{gl}_{n}(K))=n$，且 $X\in\mathfrak{gl}_{n}(K)$ 是正则的当且仅当它的所有特征值两两不同. 再证明 $\operatorname{rk}(\mathfrak{sl}_{n}(K))=n-1$.
\end{exercise*}
\begin{proposition*}
若 $X$ 是正则的，则 $\mathfrak g_{X,0}$ 是 $\mathfrak g$ 的一个 Cartan 子代数.
\end{proposition*}
\begin{proof}
i) 注意 $\mathfrak g_{X,0}$ 总是一个 Lie 子代数. 这实际上由公式（可用归纳法证明）
\[
(\operatorname{ad}_{X})^{n}([Y,Z])=\sum_{k=0}^{n}\binom{n}{k}\bigl[(\operatorname{ad}_{X})^{k}(Y),(\operatorname{ad}_{X})^{n-k}(Z)\bigr]
\]
得到.
ii) 现在设 $X$ 正则，并证明 $\mathfrak g_{X,0}$ 是幂零的. 只需证明：对一切 $Y\in\mathfrak g_{X,0}$，自同态 $(\operatorname{ad}_{Y})|_{\mathfrak g_{X,0}}$ 是幂零的. 为简化记号，令 $\mathfrak h:=\mathfrak g_{X,0}$，$r:=\operatorname{rk}(\mathfrak g)$. 我们把 $\operatorname{ad}_{Y}$ 在 $\operatorname{End}_{K}(\mathfrak g)$ 中的特征多项式记为 $\Phi^{\mathfrak g}_{\operatorname{ad}_{Y}}\in K[T]$. 由于 $Y$ 稳定 $\mathfrak h$，我们有分解 $\Phi^{\mathfrak g}_{\operatorname{ad}_{Y}}=\Phi^{\mathfrak h}_{\operatorname{ad}_{Y}}\,\Phi^{\mathfrak g/\mathfrak h}_{\operatorname{ad}_{Y}}$，而我们要证明 $\Phi^{\mathfrak h}_{\operatorname{ad}_{Y}}=T^{r}$. 取定 $\mathfrak h$ 的一组基 $X_{1},\cdots,X_{r}$，并写 $Y=\sum_{i}y_{i}X_{i}$，其中 $y_{i}\in K$. 那么存在多项式 $\Phi^{\mathfrak g}$、$\Phi^{\mathfrak h}$ 与 $\Phi^{\mathfrak g/\mathfrak h}\in K[Y_{1},\cdots,Y_{r},T]$——由于 $K$ 是无限的，它们是唯一的——经特化 $Y_{i}\mapsto y_{i}$ 后分别给出 $\Phi^{\mathfrak g}_{\operatorname{ad}_{Y}}$、$\Phi^{\mathfrak h}_{\operatorname{ad}_{Y}}$ 与 $\Phi^{\mathfrak g/\mathfrak h}_{\operatorname{ad}_{Y}}$. 在唯一分解整环 $K[Y_{1},\cdots,Y_{r},T]$ 中还有分解 $\Phi^{\mathfrak g}=\Phi^{\mathfrak h}\Phi^{\mathfrak g/\mathfrak h}$. 于是我们知道：
\begin{itemize}
\item 由 $X$ 的正则性，$T$ 在 $K[T]$ 中不整除 $\Phi^{\mathfrak g/\mathfrak h}_{\operatorname{ad}_{X}}$. 由于 $\Phi^{\mathfrak g/\mathfrak h}_{\operatorname{ad}_{X}}$ 是 $\Phi^{\mathfrak g/\mathfrak h}$ 的一个特化，故 $T$ 在 $K[Y_{1},\cdots,Y_{r},T]$ 中不整除 $\Phi^{\mathfrak g/\mathfrak h}$.
\item 由秩的定义，对一切 $Y$，$T^{r}$ 在 $K[T]$ 中整除 $\Phi^{\mathfrak g}_{\operatorname{ad}_{Y}}$. 把 $\Phi^{\mathfrak g}$ 写成 $\sum_{k}f_{k}T^{k}$，其中 $f_{k}\in K[Y_{1},\cdots,Y_{r}]$，即可推出 $f_{0}=\cdots=f_{r-1}=0$，从而 $T^{r}$ 在 $K[Y_{1},\cdots,Y_{r},T]$ 中整除 $\Phi^{\mathfrak g}$.
\end{itemize}
由此推出 $T^{r}$ 在 $K[Y_{1},\cdots,Y_{r},T]$ 中整除 $\Phi^{\mathfrak h}$，从而对一切 $Y$ 整除 $\Phi^{\mathfrak h}_{\operatorname{ad}_{Y}}$. 由于 $\Phi^{\mathfrak h}_{\operatorname{ad}_{Y}}$ 的次数为 $r$，故 $\Phi^{\mathfrak h}_{\operatorname{ad}_{Y}}=T^{r}$，如所欲证.
iii) 剩下要验证 $\mathfrak h=N(\mathfrak h)$. 但由于 $X\in\mathfrak h$，我们有 $\operatorname{ad}_{X}(N(\mathfrak h))\subset\mathfrak h$，于是 $(\operatorname{ad}_{X})^{r+1}(N(\mathfrak h))=0$，特别地 $N(\mathfrak h)\subset\mathfrak g_{X,0}=\mathfrak h$.
\end{proof}
\begin{remark*}
可以证明——但我们不打算这么做——任何 Cartan 子代数都具有这种形式.
\end{remark*}
\begin{example*}
取 $\mathfrak g=\mathfrak{sl}_{n}(K)$ 与 $X=\sum_{i=1}^{n}x_{i}E_{ii}$，其中诸 $x_{i}$ 两两不同. 那么（习题）可以算出 $\mathfrak g_{X,0}$ 是（迹为零的）对角矩阵子代数. 特别地，此时 $\mathfrak g_{X,0}$ 是一个阿贝尔的 Cartan 子代数.
\end{example*}
\begin{numpara}[Cartan 子代数的根]\label{np:lie-2-5-2}
设 $\mathfrak g$ 是一个 $K$-Lie 代数，$\mathfrak h\subset\mathfrak g$ 是一个 Cartan 子代数. 让 $\mathfrak h$ 通过 $\mathfrak g$ 的伴随表示 $(\mathfrak g,\operatorname{ad})$ 在 $\mathfrak h$ 上的限制而作用在 $\mathfrak g$ 上. 对每个线性形式 $\lambda\in\mathfrak h_{\mathrm{ab}}^{*}$，我们有广义特征空间
\[
\mathfrak g^{\lambda}=\bigl\{Y\in\mathfrak g,\ \forall X\in\mathfrak h,\ (\operatorname{ad}_{X}-\lambda(X))^{n}(Y)=0\ \text{当 }n\geqslant\dim\mathfrak g\bigr\}.
\]
\end{numpara}
\begin{lemma*}
我们有 $\mathfrak g^{0}=\mathfrak h$，并且对一切 $\lambda,\mu\in\mathfrak h_{\mathrm{ab}}^{*}$，有 $[\mathfrak g^{\lambda},\mathfrak g^{\mu}]\subset\mathfrak g^{\lambda+\mu}$.
\end{lemma*}
\begin{proof}
由于 $\mathfrak h$ 是幂零的，$(\operatorname{ad}_{X})|_{\mathfrak h}$ 对一切 $X\in\mathfrak h$ 都是幂零的，故 $\mathfrak h\subset\mathfrak g^{0}$. 反之，$\mathfrak h$ 通过幂零算子作用在 $\mathfrak g^{0}/\mathfrak h$ 上. 于是若 $\mathfrak g^{0}\supsetneq\mathfrak h$，则存在 $X\in\mathfrak g^{0}\setminus\mathfrak h$，其像 $\overline{X}$ 在 $\mathfrak g^{0}/\mathfrak h$ 中被 $\mathfrak h$ 零化. 但这意味着 $[\mathfrak h,X]\subset\mathfrak h$，故 $X\in N(\mathfrak h)=\mathfrak h$：矛盾.
第二个断言由公式
\[
(\operatorname{ad}_{X}-\lambda(X)-\mu(X))^{n}([Y,Z])=\sum_{k=0}^{n}\binom{n}{k}\bigl[(\operatorname{ad}_{X}-\lambda(X))^{k}(Y),(\operatorname{ad}_{X}-\mu(X))^{n-k}(Z)\bigr]
\]
得到，证明留给读者.
\end{proof}
\begin{definition*}
$\mathfrak h$ 在 $\mathfrak g$ 中的一个 \newterm{lie-root}{根}{root} 是 $\mathfrak h$ 在 $(\mathfrak g,\operatorname{ad})$ 中的一个非零权. 我们把根的全体记作 $\Phi$.
\end{definition*}
关于幂零代数的表示的定理给了我们一个分解
\[
\mathfrak g=\bigoplus_{\alpha\in\Phi\cup\{0\}}\mathfrak g^{\alpha},
\]
而前面的命题告诉我们：对 $\alpha,\beta\in\Phi$，括号 $[\mathfrak g^{\alpha},\mathfrak g^{\beta}]$ 在 $\alpha+\beta\notin\Phi\cup\{0\}$ 时为零，在 $\alpha+\beta=0$ 时含于 $\mathfrak h$，在 $\alpha+\beta\in\Phi$ 时含于 $\mathfrak g^{\alpha+\beta}$.
\begin{center}
\textbf{从现在起假设 $\mathfrak g$ 是半单的.}
\end{center}
我们记 $\langle X,Y\rangle=B_{\operatorname{ad}}(X,Y)$ 为 $\mathfrak g$ 的 Killing 型，因而它是非退化的. 我们的目标是证明下面的定理：
\begin{theorem}\label{thm:lie-2-5-3}
设 $\mathfrak g$ 半单，$\mathfrak h$ 为一个 Cartan 子代数. 则有：
\begin{enumerate}[label=\roman*)]
\item $\mathfrak h$ 是阿贝尔的，$\Phi$ 生成 $\mathfrak h^{*}$，且 $\langle.,.\rangle|_{\mathfrak h}$ 非退化，由公式
\[
\langle H,H'\rangle=\sum_{\alpha\in\Phi}\alpha(H)\alpha(H')
\]
给出；
\item 对一切 $\alpha\in\Phi$，$\mathfrak g^{\alpha}=\mathfrak g_{\alpha}$ 的维数为 $1$，$[\mathfrak g_{\alpha},\mathfrak g_{-\alpha}]$ 亦然；
\item $K$-向量空间 $\mathfrak{sl}_{\alpha}:=\mathfrak g_{-\alpha}\oplus[\mathfrak g_{\alpha},\mathfrak g_{-\alpha}]\oplus\mathfrak g_{\alpha}$ 是 $\mathfrak g$ 的一个 Lie 子代数，同构于 $\mathfrak{sl}_{2}(K)$. 更确切地说：
\begin{itemize}
\item 存在唯一的 $H_{\alpha}\in[\mathfrak g_{\alpha},\mathfrak g_{-\alpha}]$ 使得 $\alpha(H_{\alpha})=2$；
\item 若取定 $E_{\alpha}\in\mathfrak g_{\alpha}\setminus\{0\}$，则存在唯一的 $F_{\alpha}\in\mathfrak g_{-\alpha}$ 使得 $[E_{\alpha},F_{\alpha}]=H_{\alpha}$.
\end{itemize}
此时三元组 $(E_{\alpha},H_{\alpha},F_{\alpha})$ 是一个 \newterm{lie-sl2-triplet}{$\mathfrak{sl}_{2}$-三元组}{$\mathfrak{sl}_2$-triplet}；
\item 对 $\alpha,\beta\in\Phi$ 且 $\beta\neq-\alpha$，有 $[\mathfrak g_{\alpha},\mathfrak g_{\beta}]=\mathfrak g_{\alpha+\beta}$.
\end{enumerate}
\end{theorem}
\begin{example*}
考察 $\mathfrak g=\mathfrak{sl}_{n}(K)$ 而 $\mathfrak h$ 为对角矩阵 Cartan 子代数的情形. 公式 $\bigl[\sum x_{ii}E_{ii},\sum_{i,j}a_{ij}E_{ij}\bigr]=\sum_{i,j}(x_{ii}-x_{jj})a_{ij}E_{ij}$ 表明
\[
\Phi=\{\alpha_{ij}\in\mathfrak h^{*},\ 1\leqslant i\neq j\leqslant n\},\quad\text{其中}\ \alpha_{ij}\Bigl(\sum_{i}x_{ii}E_{ii}\Bigr)=x_{ii}-x_{jj},
\]
且 $\mathfrak g_{\alpha_{ij}}=KE_{ij}$. 此外，$[\mathfrak g_{\alpha_{ij}},\mathfrak g_{\alpha_{ji}}]=K(E_{ii}-E_{jj})$，而 $(E_{ij},E_{ii}-E_{jj},E_{ji})$ 是一个 $\mathfrak{sl}_{2}$-三元组.
\end{example*}
为了证明该定理，我们先从下面的结果开始.
\begin{proposition}\label{prop:lie-2-5-4}
设 $\mathfrak g$ 半单，$\mathfrak h$ 为一个 Cartan 子代数. 则有：
\begin{enumerate}[label=\roman*)]
\item 对 $\alpha,\beta\in\Phi\cup\{0\}$ 且 $\alpha+\beta\neq0$，有 $\langle\mathfrak g^{\alpha},\mathfrak g^{\beta}\rangle=0$；
\item $\langle.,.\rangle|_{\mathfrak h}$ 非退化；
\item $\langle.,.\rangle|_{\mathfrak g^{\alpha}\times\mathfrak g^{-\alpha}}$ 在 $\mathfrak g^{\alpha}$ 与 $\mathfrak g^{-\alpha}$ 之间给出一个对偶.
\end{enumerate}
\end{proposition}
\begin{proof}
i) 设 $\lambda\in\Phi\cup\{0\}$，$X\in\mathfrak g^{\alpha}$，$Y\in\mathfrak g^{\beta}$. 那么 $\operatorname{ad}_{X}\operatorname{ad}_{Y}(\mathfrak g^{\lambda})\subset\mathfrak g^{\lambda+\alpha+\beta}$ 且 $\lambda+\alpha+\beta\neq\lambda$. 于是在服从分解 $\mathfrak g=\bigoplus_{\lambda}\mathfrak g^{\lambda}$ 的一组基中，$\operatorname{ad}_{X}\operatorname{ad}_{Y}$ 的矩阵具有零对角线，故 $\operatorname{tr}(\operatorname{ad}_{X}\operatorname{ad}_{Y})=0$.
ii) 设 $X\in\mathfrak h$. 根据 i)，对每个根 $\alpha\in\Phi$ 都有 $X\in\mathfrak g^{\alpha,\perp}$. 于是若 $X\in\mathfrak h^{\perp}$，则 $X\in\mathfrak g^{\perp}$，从而 $X=0$.
iii) 设 $X\in\mathfrak g^{\alpha}$. 根据 i)，对每个权 $\beta\neq-\alpha$ 都有 $X\in\mathfrak g^{\beta,\perp}$. 于是若 $X\in\mathfrak g^{\perp_{-\alpha}}$，则 $X\in\mathfrak g^{\perp}$，从而 $X=0$.
\end{proof}
\begin{corollary*}
同样的假设下. 则 $\mathfrak h$ 是阿贝尔的，且 $\Phi$ 生成其对偶 $\mathfrak h^{*}$.
\end{corollary*}
\begin{proof}
i) 设 $\mathcal V$ 是 $\mathfrak g$ 的一个旗使得 $\operatorname{ad}(\mathfrak h)\subset\mathfrak b(\mathcal V)$（因为 $\mathfrak h$ 可解）. 于是 $\operatorname{ad}([\mathfrak h,\mathfrak h])\subset\mathfrak b^{\mathrm{nilp}}(\mathcal V)$. 设 $Y\in[\mathfrak h,\mathfrak h]$. 对一切 $X\in\mathfrak h$，有 $\operatorname{ad}_{X}\operatorname{ad}_{Y}\in\mathfrak b^{\mathrm{nilp}}(\mathcal V)$，故 $\operatorname{tr}(\operatorname{ad}_{X}\operatorname{ad}_{Y})=0$. 于是 $Y\in\mathfrak h^{\perp}$，由前一命题之 ii) 得 $Y=0$. 故 $[\mathfrak h,\mathfrak h]=0$，$\mathfrak h$ 确实是阿贝尔的.
ii) 设 $Y$ 满足对一切 $\alpha\in\Phi$ 有 $\alpha(Y)=0$. 那么 $\operatorname{ad}_{Y}$ 的所有特征值都为零，故 $\operatorname{ad}_{Y}$ 是幂零的. 与上面一样可证 $Y\in\mathfrak h^{\perp}$，而该命题之 ii) 给出 $Y=0$.
\end{proof}
\begin{proposition}\label{prop:lie-2-5-5}
对每个根 $\alpha\in\Phi$，有 $\dim(\mathfrak g^{\alpha})=1$（从而 $\mathfrak g^{\alpha}=\mathfrak g_{\alpha}$），且 $\{n\in\mathbb Z,\ n\alpha\in\Phi\}=\{\pm1\}$.
\end{proposition}
\begin{proof}
设 $E_{\alpha}\in\mathfrak g_{\alpha}$（注意，$\mathfrak g_{\alpha}$ 是特征空间，由于 $\alpha$ 是一个权，它非零）. 于是对一切 $H\in\mathfrak h$ 有 $[H,E_{\alpha}]=\alpha(H)E_{\alpha}$. 选取 $F\in\mathfrak g^{-\alpha}$ 使得 $\langle E_{\alpha},F\rangle=1$，并令 $H'_{\alpha}:=[E_{\alpha},F]\in\mathfrak h=\mathfrak g^{0}$. 那么对一切 $H\in\mathfrak h$，有
\[
(\ast)\qquad\langle H,H'_{\alpha}\rangle=\langle H,[E_{\alpha},F]\rangle=\langle[H,E_{\alpha}],F\rangle=\alpha(H)\langle E_{\alpha},F\rangle=\alpha(H).
\]
我们来证明 $\alpha(H'_{\alpha})\neq0$. 实际上，上面的推论保证了存在 $\beta$ 使得 $\beta(H'_{\alpha})\neq0$. 考虑 $W:=\sum_{n\in\mathbb Z}\mathfrak g^{\beta+n\alpha}$. 那么 $W$ 在 $\operatorname{ad}_{E_{\alpha}}$ 与 $\operatorname{ad}_{F}$ 下稳定，故也在 $\operatorname{ad}_{[E_{\alpha},F]}=\operatorname{ad}_{H'_{\alpha}}$ 下稳定，并且还有 $\operatorname{tr}((\operatorname{ad}_{H'_{\alpha}})|_{W})=\operatorname{tr}([( \operatorname{ad}_{E_{\alpha}})|_{W},(\operatorname{ad}_{F})|_{W}])=0$. 而 $\operatorname{tr}((\operatorname{ad}_{H'_{\alpha}})|_{W})=\sum_{n\in\mathbb Z}\operatorname{tr}((\operatorname{ad}_{H'_{\alpha}})|_{\mathfrak g^{\beta+n\alpha}})=\sum_{n\in\mathbb Z}(\beta+n\alpha)(H'_{\alpha})\dim(\mathfrak g^{\beta+n\alpha})$.\footnote{译注：原文此处最后一个 $\dim$ 带下标 $\mathfrak g^{\beta+n\alpha}$，按上下文应为 $\dim(\mathfrak g^{\beta+n\alpha})$.} 由于 $\beta(H'_{\alpha})\neq0$，必须有 $\alpha(H'_{\alpha})\neq0$.
我们来证明 $\operatorname{ad}_{E_{\alpha}}(\mathfrak g^{-\alpha})\subset K.H'_{\alpha}$. 实际上，对一切 $Y\in\mathfrak g^{-\alpha}$ 与 $H\in\mathfrak h$，有 $\langle H,[E_{\alpha},Y]\rangle=\langle[H,E_{\alpha}],Y\rangle=\alpha(H)\langle E_{\alpha},Y\rangle$，故由 $(\ast)$ 得 $[E_{\alpha},Y]=\langle E_{\alpha},Y\rangle H'_{\alpha}$.
现在令
\[
V:=K.E_{\alpha}\oplus K.H'_{\alpha}\oplus\Bigl(\bigoplus_{n<0}\mathfrak g^{n\alpha}\Bigr).
\]
由构造，$V$ 在 $\operatorname{ad}_{F}$ 下稳定，而我们刚刚看到它也在 $\operatorname{ad}_{E_{\alpha}}$ 下稳定. 与上面一样，由此推出它在 $\operatorname{ad}_{H'_{\alpha}}$ 下稳定，且 $\operatorname{tr}((\operatorname{ad}_{H'_{\alpha}})|_{V})=0$. 而 $\operatorname{tr}((\operatorname{ad}_{H'_{\alpha}})|_{V})=\alpha(H'_{\alpha})\bigl(1+\sum_{n<0}n\dim(\mathfrak g^{n\alpha})\bigr)$. 由于 $\alpha(H'_{\alpha})\neq0$，定理得证.
\end{proof}
\begin{corollary*}
对 $H,H'\in\mathfrak h$，有 $\langle H,H'\rangle=\sum_{\alpha\in\Phi}\alpha(H)\alpha(H')$.
\end{corollary*}
\begin{proof}
显然，因为 $\operatorname{ad}_{H}\operatorname{ad}_{H'}(\mathfrak g^{\alpha})\subset\mathfrak g^{\alpha}$ 且 $\mathfrak g^{\alpha}=\mathfrak g_{\alpha}$ 的维数为 $1$.
\end{proof}
\begin{corollary*}
定理~\ref{thm:lie-2-5-3} 之点 iii) 成立.
\end{corollary*}
\begin{proof}
沿用前一证明的记号，我们有 $\operatorname{ad}_{E_{\alpha}}(\mathfrak g_{-\alpha})=K.H'_{\alpha}$，故 $[\mathfrak g_{\alpha},\mathfrak g_{-\alpha}]=K.H'_{\alpha}$ 的维数为 $1$. 于是存在唯一的 $H_{\alpha}\in[\mathfrak g_{\alpha},\mathfrak g_{-\alpha}]$ 使得 $\alpha(H_{\alpha})=2$. 具体地，有 $H_{\alpha}=\dfrac{2H'_{\alpha}}{\alpha(H'_{\alpha})}=\dfrac{2H'_{\alpha}}{\langle H'_{\alpha},H'_{\alpha}\rangle}$. 现在存在唯一的 $F_{\alpha}\in\mathfrak g_{-\alpha}$ 使得 $\langle E_{\alpha},F_{\alpha}\rangle=\dfrac{2}{\langle H'_{\alpha},H'_{\alpha}\rangle}$. 这时 $[E_{\alpha},F_{\alpha}]=H_{\alpha}$，从而三元组 $(E_{\alpha},H_{\alpha},F_{\alpha})$ 是一个 $\mathfrak{sl}_{2}$-三元组.
\end{proof}
\begin{proposition}\label{prop:lie-2-5-6}
设 $\beta\in\Phi$ 满足 $\beta\neq\pm\alpha$. 则：
\begin{enumerate}[label=\roman*)]
\item $\beta(H_{\alpha})\in\mathbb Z$，且 $\beta-\beta(H_{\alpha})\in\Phi$；\footnote{译注：原文如此；由其证明末尾的``$\mathfrak g_{\beta-\beta(H_{\alpha})\alpha}\neq0$''可知应为 $\beta-\beta(H_{\alpha})\alpha\in\Phi$.}
此外，$\{n\in\mathbb Z,\ \beta+n\alpha\in\Phi\}$ 是 $\mathbb Z$ 中的一个区间.
\item 若 $\beta+\alpha\in\Phi$，则 $[\mathfrak g_{\alpha},\mathfrak g_{\beta}]=\mathfrak g_{\alpha+\beta}$.
\end{enumerate}
\end{proposition}
\begin{proof}
考虑 $\mathfrak g_{\beta+\mathbb Z\alpha}:=\bigoplus_{n\in\mathbb Z}\mathfrak g_{\beta+n\alpha}$. 它是 $\mathfrak g$ 的一个在 $\operatorname{ad}_{E_{\alpha}}$ 与 $\operatorname{ad}_{F_{\alpha}}$ 下稳定的子空间，因而是 $\mathfrak{sl}_{\alpha}$ 的一个表示. $H_{\alpha}$ 在 $\mathfrak g_{\beta+\mathbb Z\alpha}$ 中的权为\footnote{译注：下式中 $\mathfrak g_{\alpha+n\beta}$ 原文如此，按上下文应为 $\mathfrak g_{\beta+n\alpha}$.}
\[
\{\beta(H_{\alpha})+2n,\ \text{其中 }n\ \text{满足 }\ \mathfrak g_{\alpha+n\beta}\neq0\}.
\]
$\mathfrak{sl}_{2}(K)$ 的表示论告诉我们这些权都是整数，故 $\beta(H_{\alpha})\in\mathbb Z$. 它还告诉我们，由于这些权的重数均为 $1$ 且具有相同的奇偶性，$\mathfrak g_{\beta+\mathbb Z\alpha}$ 是 $\mathfrak{sl}_{\alpha}$ 的一个不可约表示，且它的权也形如 $\{-m,-m+2,\cdots,m-2,m\}$. 由此我们得到 i) 中的区间性质，并且由于 $\beta(H_{\alpha})$ 是一个权，$-\beta(H_{\alpha})$ 也是一个权. 于是 $\mathfrak g_{\beta-\beta(H_{\alpha})\alpha}\neq0$，从而 $\beta-\beta(H_{\alpha})\alpha\in\Phi$.
ii) 设 $\mathfrak g_{\beta+\alpha}\neq0$. 它是 $H_{\alpha}$ 的权为 $\beta(H_{\alpha})+2$ 的权空间. 根据 $\mathfrak{sl}_{2}$ 不可约表示的已知结构，有 $[\mathfrak g_{\alpha},\mathfrak g_{\beta}]=\operatorname{ad}_{E_{\alpha}}(\mathfrak g_{\beta})=\mathfrak g_{\alpha+\beta}$，因为 $\mathfrak g_{\beta}$ 是 $H_{\alpha}$ 的权为 $\beta(H_{\alpha})$ 的权空间.
\end{proof}
至此我们完成了定理~\ref{thm:lie-2-5-3} 的证明. 上面的命题还有下述推论：
\begin{corollary*}
对 $\alpha\in\Phi$，有 $\Phi\cap K\alpha=\{\pm\alpha\}$.
\end{corollary*}
\begin{proof}
假设存在 $\beta\in\Phi\cap K\alpha$，它异于 $\alpha$ 与 $-\alpha$. 由于 $\beta(H_{\alpha})\in\mathbb Z$，有 $\beta\in\frac{1}{2}\mathbb Z\alpha$. 写 $\beta=\frac{n}{2}\alpha$，其中 $n$ 为奇数. 那么由于 $-\beta\in\Phi$，该命题之 i) 的区间性质保证 $\frac{1}{2}\alpha\in\Phi$. 但根据命题~\ref{prop:lie-2-5-5}，$\Phi$ 中 $\frac{1}{2}\alpha$ 的倍数只有 $\pm\frac{1}{2}\alpha$，这与 $\alpha\in\Phi$ 矛盾.
\end{proof}
\begin{numpara}[$\Phi$ 的有理性]\label{np:lie-2-5-7}
我们把迄今为止关于集合 $\Phi$ 所证明的结果概括如下：
\begin{enumerate}[label=\roman*)]
\item $\Phi$ 生成 $K$-向量空间 $\mathfrak h^{*}$；
\item 对每个根 $\alpha\in\Phi$，有 $\Phi\cap K\alpha=\{\pm\alpha\}$；
\item 对两个根 $\beta,\alpha$，有 $\beta(H_{\alpha})\in\mathbb Z$，且当 $\beta\neq-\alpha$ 时，$\{n\in\mathbb Z,\ \beta+n\alpha\in\Phi\}$ 是一个包含 $-\beta(H_{\alpha})$ 的区间.
\end{enumerate}
现在令
\[
\mathfrak h_{\mathbb Q}:=\operatorname{Vect}_{\mathbb Q}\{H_{\alpha},\ \alpha\in\Phi\},\qquad\text{及}\qquad\mathfrak h^{*}_{\mathbb Q}:=\operatorname{Vect}_{\mathbb Q}\{\alpha,\ \alpha\in\Phi\}.
\]
\end{numpara}
\begin{proposition*}
沿用上面的记号，则：
\begin{enumerate}[label=\roman*)]
\item $\dim_{\mathbb Q}(\mathfrak h_{\mathbb Q})=\dim_{\mathbb Q}(\mathfrak h^{*}_{\mathbb Q})=\operatorname{rk}(\mathfrak g)$（$=\dim_{K}(\mathfrak h)$）；
\item 对偶 $\mathfrak h\times\mathfrak h^{*}\longrightarrow K$ 诱导出一个对偶 $\mathfrak h_{\mathbb Q}\times\mathfrak h^{*}_{\mathbb Q}\longrightarrow\mathbb Q$；
\item Killing 型诱导一个有理形式 $\langle.,.\rangle:\mathfrak h_{\mathbb Q}\times\mathfrak h_{\mathbb Q}\longrightarrow\mathbb Q$；
\item $\langle.,.\rangle$ 到 $\mathfrak h_{\mathbb R}$ 上的扩张是正定的.
\end{enumerate}
\end{proposition*}
\begin{proof}
i) 设 $\alpha_{1},\cdots,\alpha_{r}$ 是由 $\Phi$ 中元素组成的 $\mathfrak h^{*}$ 的一组基. 那么诸 $H_{\alpha_{i}}$ 构成 $\mathfrak h$ 的一组 $K$-基（注意 $\langle H_{\alpha_{i}},H\rangle=\frac{\langle H_{\alpha_{i}},H_{\alpha_{i}}\rangle}{2}\,\alpha_{i}(H)=\frac{2}{\langle H'_{\alpha_{i}},H'_{\alpha_{i}}\rangle}\,\alpha_{i}(H)$）. 对 $\alpha\in\Phi$，写 $H_{\alpha}=\sum_{i=1}^{r}\lambda_{i}H_{\alpha_{i}}$，其中 $\lambda_{i}\in K$. 我们必须证明 $\lambda_{i}\in\mathbb Q$，$\forall i$. 设 $M\in\mathrm M_{r}(K)$ 是矩阵 $\bigl(\langle H_{\alpha_{i}},H_{\alpha_{j}}\rangle\bigr)_{1\leqslant i,j\leqslant r}$. 那么 $M$ 是可逆的（因为 $\langle.,.\rangle$ 非退化），而 $\lambda=(\lambda_{1},\cdots,\lambda_{r})\in K^{r}$ 是线性方程组 $M.\lambda=\mu_{\alpha}$ 的唯一解，其中 $\mu_{\alpha}=\bigl(\langle H_{\alpha_{1}},H_{\alpha}\rangle,\cdots,\langle H_{\alpha_{r}},H_{\alpha}\rangle\bigr)\in K^{r}$. 而对两个根 $\alpha,\beta$，有
\[
\langle H_{\alpha},H_{\beta}\rangle=\sum_{\gamma\in\Phi}\gamma(H_{\alpha})\gamma(H_{\beta})\in\mathbb Z,
\]
故 $M\in\mathrm M_{r}(\mathbb Q)$，$\mu_{\alpha}\in\mathbb Q^{r}$，最终 $\lambda\in\mathbb Q^{r}$. 由此推出诸 $H_{\alpha_{i}}$（$i=1,\cdots,r$）构成 $\mathfrak h_{\mathbb Q}$ 的一组 $\mathbb Q$-基. 用同样的方法可证 $\alpha_{1},\cdots,\alpha_{r}$ 构成 $\mathfrak h^{*}_{\mathbb Q}$ 的一组 $\mathbb Q$-基.
ii) $\mathfrak h$ 与 $\mathfrak h^{*}$ 之间的 $K$-对偶确实把 $\mathfrak h_{\mathbb Q}\times\mathfrak h^{*}_{\mathbb Q}$ 映入 $\mathbb Q$，因为对一切 $\alpha,\beta\in\Phi$ 都有 $\beta(H_{\alpha})\in\mathbb Z$. 由于 $\mathfrak h_{\mathbb Q}$（相应地，$\mathfrak h^{*}_{\mathbb Q}$）在 $K$ 上生成 $\mathfrak h$（相应地，$\mathfrak h^{*}$），该配对是完美的.
iii) 同理. 我们已经看到 $\langle\mathfrak h_{\mathbb Q},\mathfrak h_{\mathbb Q}\rangle\subset\mathbb Q$，且由于 $\mathfrak h_{\mathbb Q}$ 生成 $\mathfrak h$，该配对是非退化的.
iv) 由于 $\langle H,H\rangle=\sum_{\alpha}\alpha(H)^{2}$，所考虑的实双线性型是正的. 而我们知道它是非退化的，故它也是正定的.
\end{proof}
